%% kapa: exam=1 round=19 year=2008 period=2 form=A subject=economics subject_ko=경제원론 seq=31 number=31
%% source: Q-Net 감정평가사 기출문제 (cst00309) artlSeq=1001962 file=extracted/1008/1001962/2008년 제19회 감정평가사 1차시험 기출문제/2교시.hwp
%% conversion: kordoc 4.24.51 -> tools/gen1.py (auto); review=proofread-sonnet; flags=none
\begin{question}{31}
화장품 시장에 새로 진입을 고려 중인 B사와 이미 영업 중인 A사가 어떤 전략을 택하는가에 따라 두 기업의 보수, 즉 (A의 보수, B의 보수)는 다음과 같이 결정된다. 보수를 극대화하려는 두 기업이 취할 전략에 대한 설명으로 옳지 않은 것은?

{\small\centering
\begin{tabular}{|c|c|c|}
\hline
\makecell[l]{\hspace*{2em}B\\A} & 진입 & 진입하지 않음 \\
\hline
강한 대응 & (2, 5) & (10, 8) \\
\hline
약한 대응 & (3, 11) & (12, 8) \\
\hline
\end{tabular}\par}

\begin{choices}
  \item B사가 진입할 경우 A사는 약한 대응을 하는 것이 이득이다.
  \item A사가 약한 대응을 하고, B사가 진입하는 것은 내쉬균형이다.
  \item A사는 약한 대응을 하는 것이 우월 전략이다.
  \item A사가 약한 대응을 할 것으로 예상되면 B사는 진입하는 것이 이득이다.
  \item ‘B사가 진입하면 강한 대응을 하겠다’라는 A사의 경고는 신뢰할 만하다.
\end{choices}
\end{question}
\answer{5}
